% \chapter{Proposed Approach}

% Overview of Your Proposed Framework: Briefly introduce your idea for a framework (even though it’s early, you can outline its broad components and goals). This should focus on the methodology you plan to use, such as AI model development, data collection, and integration techniques.
% Key Challenges in Building the Framework: Highlight some of the challenges you foresee in building this framework, such as data quality, model robustness, or computational efficiency.
% Potential Impact: Discuss how your framework, once developed, could advance the field of motion tracking or be applied to specific industries.

\chapter{Methodology}

% Research Approach: Outline the approach you plan to take for your research in the next semester. For example, will you focus on developing specific AI models, creating a tool for motion tracking, or synthesizing a new algorithm?
% Data Collection Strategy: Explain your plan for gathering datasets or experimental data (e.g., existing motion tracking datasets or data collection methods).
% Evaluation Metrics: Define what metrics will be important for evaluating the effectiveness of the framework once developed (e.g., accuracy, processing speed, real-time tracking capability).

\section{Overview}
% The research methodology for this study combines theoretical exploration, empirical validation, and practical implementation to develop a robust framework for generating AI models tailored to motion tracking applications. This chapter details the structured approach employed to achieve the research objectives while ensuring replicability and relevance to real-world scenarios. By adopting a modular and iterative design process, the study addresses critical challenges such as resource constraints on edge devices, data preprocessing complexities, and the need for a user-friendly interface. The methodology ensures the alignment of each research activity with the overarching goals and research questions.

% The sections below explore the research design, tools and technologies used, data collection strategies, preprocessing methods, framework development process, evaluation metrics, and ethical considerations.

The proposed framework this study is designed to develop robust and efficient AI models for motion tracking applications, focusing on practical implementation, scalability, and ease of use. Combining theoretical exploration, empirical validation, and modular design, the framework addresses key challenges such as data quality, computational constraints, and user accessibility. By integrating data preprocessing, model development, and a user-friendly interface, it aims to deliver a comprehensive and effective solution for motion tracking systems.\\

The framework is structured around three key components:

\begin{itemize}
    \item \textbf{Data Collection and Preprocessing}
    
    Motion data collected from various sensors, such as accelerometers and gyroscopes, often produces noisy and inconsistent outputs. To ensure reliable data and model training readiness, the framework uses robust preprocessing techniques, including filtering, normalization, and data augmentation. These methods help reduce noise, handle missing values, and enrich the dataset, improving the generalizability of models.

    \item \textbf{AI Model Development and Optimization}
    
    The framework emphasizes the creation of lightweight and efficient AI architectures specifically designed for resource-constrained environments such as edge devices. The models are designed to achieve high accuracy and real-time performance while maintaining computational efficiency. The iterative development process ensures continuous refinement, adapting models to different motion patterns and tracking requirements.

    \item \textbf{System Integration and User Interface}

    Seamless integration across multiple hardware and software environments is a core focus of the framework. Additionally, a user-friendly interface allows users to visualize motion data, interpret model output, and customize system settings with ease. This design prioritizes accessibility, allowing both expert and non-expert users to use the system effectively.
\end{itemize}


\textbf{KEY CHALLENGES}

Despite its well-structured design, the framework must address several challenges to ensure success:
\begin{itemize}
    \item \textbf{Data quality and variability}
    
    Motion sensor data often contains noise, inconsistencies, and missing values, which can hinder model performance. Advanced preprocessing techniques are required to normalize and prepare the data for training.

    \item \textbf{Integration and usability}
    
    Ensuring compatibility with a variety of hardware setups and maintaining a user-friendly experience for non-experts requires careful design and testing.

    \item \textbf{Privacy and ethical concerns}
    
    The collection of motion data raises privacy concerns. The framework integrates data anonymization methods and compliance with data protection regulations to ensure the ethical use of data.
\end{itemize}


\textbf{POTENTIAL IMPACTS}

The proposed framework has the potential to bring transformative benefits to the motion tracking field and its applications:

\begin{itemize}
    \item \textbf{Advanced motion tracking technology}
    
    By providing a robust, scalable, and efficient solution, the framework can significantly improve the accuracy and reliability of motion tracking systems, enabling detailed analysis of human movements.

    \item \textbf{Enabling industry applications}

    From healthcare (e.g., rehabilitation and physical therapy) to fitness (e.g., activity tracking) and entertainment (e.g., motion capture in games and movies), the framework enables versatile applications across multiple domains.

    \item \textbf{Democratizing Accessibility}
    
    A focus on user-friendly design ensures that the technology is accessible to a wider audience, including researchers, developers, and practitioners with limited technical expertise.
\end{itemize}

By addressing these challenges and emphasizing innovation, scalability, and usability, the proposed framework aims to set a new standard for motion tracking systems. Its potential to revolutionize motion tracking technology opens up exciting possibilities for advancing research and industry applications.

\section{Research Design}
The research design outlines the high-level structure of the study, defining how the objectives will be translated into actionable steps. It emphasizes a systematic, iterative approach to ensure flexibility and adaptability during the research process. The design is divided into the following stages:

\subsection{Literature Review}
The first stage involves an in-depth examination of existing literature in motion tracking systems, AI model generation frameworks, and edge computing technologies. Key objectives include:

\begin{itemize}
    \item Identifying current advancements in motion tracking and their application areas (e.g., healthcare, sports, and human-computer interaction).
    \item Analyzing the capabilities and limitations of frameworks like Edge Impulse and Google MediaPipe.
    \item Understanding the challenges posed by motion data, such as noise, heterogeneity, and the need for accurate annotation.
\end{itemize}
This stage sets the foundation for identifying gaps and opportunities that inform the design of the proposed framework.

\subsection{Requirement Analysis}
Requirement analysis bridges the gap between theoretical insights and practical implementation. Key tasks include:

Functional Requirements:
\begin{itemize}
    \item Support for multiple input sources, including IMU sensors and external datasets.
    \item Capability to automate data preprocessing, feature extraction, and model training.
    \item Seamless integration with edge devices and mobile applications.
\end{itemize}
Non-Functional Requirements:
\begin{itemize}
    \item Scalability to handle increasing data volumes or new sensor modalities.
    \item Resource efficiency, ensuring models are lightweight and computationally feasible for edge devices.
    \item User-friendliness to cater to both technical and non-technical users.
\end{itemize}
These requirements serve as benchmarks during the framework development and evaluation stages.

\subsection{Framework Development}
The framework is developed iteratively to incorporate user feedback and address technical challenges. This process involves:

\begin{itemize}
    \item Designing a modular architecture that separates preprocessing, model generation, and deployment components.
    \item Leveraging existing frameworks, such as TensorFlow Lite, for rapid prototyping while ensuring the flexibility to integrate custom algorithms.
    \item Incorporating a graphical user interface (GUI) to simplify user interactions.
\end{itemize}

\subsection{Experimental Validation}
Validation ensures that the framework meets its functional and performance requirements. Experiments are conducted to evaluate the framework under diverse conditions:

\begin{itemize}
    \item Data Variability: Test with various motion types and intensities.
    \item Hardware Constraints: Assess performance on edge devices like the Seeed XIAO nRF52840 Sense.
    \item User Scenarios: Simulate real-world applications, such as fitness tracking or gesture recognition.
\end{itemize}
\subsection{Evaluation and Iteration}
Evaluation is an ongoing process, with feedback loops integrated at every stage. Metrics such as accuracy, latency, and user satisfaction guide refinements to the framework. Multiple iterations are conducted to ensure robustness and usability.

\section{Tools and Technologies}
The choice of tools and technologies is critical for addressing the unique requirements of motion tracking applications. A combination of hardware, software, and development platforms is used to create and validate the framework.

\subsection{Hardware}
\begin{itemize}
    \item Seeed XIAO nRF52840 Sense:
    This microcontroller serves as the primary sensor hub, equipped with an onboard IMU to capture motion data and BLE capabilities for wireless communication. Its compact size and low power consumption make it ideal for edge applications.
    \item BLE-Enabled Smartphones:
    Smartphones are used to receive data transmitted by the Seeed XIAO and provide a user-friendly interface for data visualization and analysis.
\end{itemize}

\subsection{Software and Libraries}
\begin{itemize}
    \item Edge Impulse:
    A no-code AI development platform, Edge Impulse enables rapid prototyping of AI models for edge devices. Its preprocessing and feature extraction pipelines are particularly relevant for motion data.
    \item Google MediaPipe:
    MediaPipe's gesture recognition pipelines offer insights into designing efficient, real-time motion tracking models.
    \item TensorFlow Lite:
    TensorFlow Lite allows for the deployment of lightweight AI models on edge devices, optimizing for speed and memory efficiency.
    \item Arduino IDE and PlatformIO:
    These platforms are used to program the Seeed XIAO and develop firmware for sensor data collection and transmission.
\end{itemize}

\subsection{Development Environment}
A hybrid environment is adopted, combining:

\begin{itemize}
    \item Local Tools: Python scripts for preprocessing and model training.
    \item Cloud Platforms: Edge Impulse Studio for data preprocessing and AI model benchmarking.
\end{itemize}

\section{Data Collection and Preprocessing}
\subsection{Data Sources}
Primary Data:
Raw motion data collected using the IMU sensors on the Seeed XIAO nRF52840 Sense.
Supplementary Data:
Synthetic datasets generated to complement real-world data, ensuring diverse motion scenarios.
\subsection{Preprocessing Steps}
\begin{itemize}
    \item Noise Reduction:
     Filtering techniques, such as low-pass or Kalman filters, are applied to mitigate sensor noise.
    \item Segmentation:
    Motion data is divided into fixed-size windows, enabling the identification of discrete activities.
    \item Feature Extraction:
    Key features are extracted from each segment, including:
    \begin{itemize}
        \item Statistical metrics (e.g., mean, variance).
        \item Frequency-domain features (e.g., spectral energy).
        \item Derived metrics like acceleration magnitude and orientation angles.
    \end{itemize}
    \item Data Normalization:
    All features are scaled to a uniform range for consistency during model training.
\end{itemize}

\section{Framework Development}
\subsection{Architecture}
The framework is designed to operate as an end-to-end system, consisting of:
\begin{itemize}
    \item Input Module:
    Handles raw data ingestion from IMU sensors.
    \item Processing Module:
    Automates noise filtering, segmentation, and feature extraction.
    \item Modeling Module:
    Supports training of AI models using predefined templates or custom algorithms.
    \item Output Module:
    Generates optimized models for deployment on edge devices.
\end{itemize}
\subsection{Workflow}
\begin{enumerate}
    \item Data is collected from the Seeed XIAO and ingested into the framework.
    \item Preprocessing pipelines clean and prepare the data.
    \item AI models are trained, evaluated, and exported for deployment.
\end{enumerate}

\section{Ethical Considerations}
Ethical issues are considered throughout the research, ensuring:

\begin{itemize}
    \item Data Privacy: Safeguards are in place to anonymize motion data.
    \item Inclusivity: Diverse datasets are collected to avoid algorithmic biases.
\end{itemize}
\section{Summary}
This chapter provided a comprehensive methodology for developing, implementing, and evaluating the proposed framework. Subsequent chapters will delve into the experimental setup, results, and analysis.